\begin{abstract}
How should auxiliary data be mixed when target risk has more than one valley? We study this question in a dual-domain full-ridge model with task mismatch. One coordinate cannot produce multiple valleys, while two coordinates suffice: a certified example has two local minima, and a one-parameter family exhibits valley creation, annihilation, and two distinct global minimizers. To identify all global candidates, we reduce the population stationary equation to polynomial root isolation and compare the stationary and boundary risks. Zero mismatch provides a contrasting log-convex or flat regime. We then connect this population geometry to finite samples through a fixed-dimension expansion of the actual risk and its first two derivatives. For sufficiently large samples, a nondegenerate interior optimum admits a first-order weight correction, while two nondegenerate interior population minima are selected by their unequal $R_1$ values. Under matched spectra and zero mismatch, an exact quadratic $R_1$ determines the limiting weight when its coefficients are nonzero. Validation combines exact population certificates with conditional Monte Carlo for risk and weight expansions; a fixed-candidate comparison is distinguished from finite-sample branch validation. A general finite-sample shape classification and robot-policy validation remain open.
\end{abstract}
